\chapter{Introduction to Networks}

\section{Strong and Weak Ties}
\begin{itemize}
  \item Strong ties: frequent interaction, trust, emotional closeness
  \item Weak ties: infrequent interaction, limited intimacy
  \item Strong ties: overlapping social circles
  \item Weak ties: bridges between disconnected groups
  \item Weak ties: novel information and opportunities
\end{itemize}

\section{Bridges and Local Bridges}
\begin{itemize}
  \item Bridge: an edge whose removal disconnects the network into separate components
  \item Local bridge: an edge whose endpoints do not have any direct ties to the same neighboring nodes
  \item Local bridge: removing it increases the distance between its endpoints, giving it a span of at least 3
  \item Span of a bridge: the distance between the endpoints after the edge is removed
  \item Bridges and local bridges: often connect otherwise distant or disconnected regions of the network
\end{itemize}

\section{Clustering}
\begin{itemize}
  \item Clustering: connected neighbors also connected to one another
  \item Local clustering coefficient: connected neighbor pairs divided by possible neighbor pairs
  \item Formula: $C_i = \frac{2e_i}{k_i(k_i-1)}$
  \item $e_i$: edges among node $i$'s neighbors
  \item Network clustering coefficient: average local coefficient
\end{itemize}

\section{Strong Triadic Closure Property}
\begin{itemize}
  \item Two strong ties sharing a node: likely third edge
  \item Strong ties around a node: mutually connected neighbors
  \item Open triad: two strong ties without a connecting edge
  \item Open triads: pressure toward closure
  \item Closed triads: higher clustering coefficients
\end{itemize}

\section{Homophily}
\begin{itemize}
  \item Homophily: the tendency for people to form ties with similar others
  \item Similarity may involve age, gender, education, location, interests, or beliefs
  \item Homophily can produce clustered groups with shared characteristics
  \item Network ties can reinforce existing similarities through social influence
  \item Observed similarity may result from selection, influence, or both
  \item Let $p$ and $q$ be the proportions of two groups, with $p+q=1$
  \item Under random mixing, the probability of a cross-group tie is $2pq$
  \item Cross-group tie probability $c$ under random mixing: expected to be $2pq$
  \item Homophily where the observed cross-group tie probability is less than expected: $c < 2pq$
\end{itemize}

\section{Affiliation Networks}
\begin{itemize}
  \item Affiliation network: a network connecting people to the groups, events, organizations, or activities they participate in
  \item Two-mode network: contains two types of nodes, such as individuals and affiliations
  \item One-mode projection: connects two people when they share an affiliation
  \item Shared affiliations can create indirect social ties and opportunities for interaction
  \item Affiliations can reveal overlapping communities and patterns of participation
\end{itemize}

\section{Social Affiliation Networks}
\begin{itemize}
  \item Social affiliation network: represents relationships between people and the groups, events, or organizations with which they are affiliated
  \item Membership links connect each person to the affiliations they share
  \item Shared affiliations can create indirect opportunities for social interaction and information flow
  \item One-mode projection: connects people who share at least one affiliation
  \item Projection ties may be weighted by the number of affiliations two people share
  \item Affiliation networks can reveal overlapping communities, central participants, and patterns of group involvement
\end{itemize}

\subsection{Closure in Social Affiliation Networks}

\begin{itemize}
  \item Triadic: nodes with a common neighbor tend to form a direct connection, closing the triad
  \item Group/focal: nodes connected to the same group or affiliation tend to form direct connections, closing the group-based triad
  \item Membership: nodes are more likely to join groups that their neighbors are already members of
\end{itemize}

\section{Signed Networks}

\begin{itemize}
  \item Signed network: a network in which edges have positive or negative signs, representing friendly or antagonistic relationships
  \item Positive tie: indicates friendship, trust, or alliance
  \item Negative tie: indicates enmity, distrust, or conflict
  \item Balance theory: suggests that triads tend to evolve toward a balanced state, where the product of the signs of the edges is positive
  \item Balanced triads: all positive edges or two negative edges and one positive edge
  \item Unbalanced triads: one negative edge and two positive edges or all negative edges
\end{itemize}